\chapter{Methodology and Design}
\label{chapter:methodology_and_design}

\section{Proposed Analytical Framework and Resolution Strategy}

To address the limitations identified in the current state-of-the-art—specifically computational latency, architectural 
rigidity, and the difficulty of integrating heterogeneous data—this research proposes a multi-dimensional analytical 
framework. The resolution strategy is built upon three fundamental pillars: the formalization of a multilayered discrete 
space, the adoption of a "Data-as-Interface" philosophy, and the implementation of a decoupled software architecture.


\subsection{The Multilayer Conceptualization: $m:n-CA^k$}

The core of the proposed resolution lies in the adoption of the Multilayer Cellular Automata ($m:n-CA^k$) paradigm. Unlike 
standard Cellular Automata, which operate on a single state per cell, the $m:n-CA^k$ framework allows for the simultaneous 
representation of diverse physical and environmental variables across different computational planes.

In this framework, the river basin is modeled as a stack of $n$ information layers (e.g., topography, soil saturation, 
precipitation, and urban obstacles). This allows the system to handle the complexity of flash floods by defining interaction 
rules between layers (inter-layer) as well as within the same layer (intra-layer). By treating these layers as channels of a 
global tensor, the model can execute complex transitions—such as the infiltration of water from the surface layer to the soil 
layer—without losing the computational efficiency inherent to discrete models.


\subsection{The "Data-as-Interface" Strategy}

A critical component of the resolution strategy is the "Data-as-Interface" (DaI) approach. This methodology posits that the 
data structure itself should act as the contract between different system components. Following the principles of Society 5.0 
and Open Models \parencite{FonsecaiCasas2026}, the proposed framework separates the formal specification of the model 
(what the water does) from its technical implementation (how it is programmed).

By using the IDRISI format and Hierarchical IDRISI Format (HIF), the framework ensures that both static GIS data and 
dynamic real-time sensor inputs are ingested through a standardized pipeline. This solves the "data silo" problem and allows 
the simulation engine to remain agnostic to the source of the information, whether it comes from historical databases or live 
meteorological radar.


\subsection{Architectural Resolution: Hexagonal Design}

To solve the architectural rigidity found in monolithic simulation tools, this thesis proposes a Hexagonal Architecture (also 
known as Ports and Adapters). The resolution strategy involves isolating the "Simulator Core"—which contains the $m:n-CA^k$ 
transition logic—from external dependencies such as file systems, GIS libraries, and user interfaces.

\begin{itemize}
    \item \textbf{The Core}: Encapsulates the mathematical and physical logic of the cellular automata.
    \item \textbf{The Ports}: Define the required inputs (e.g., rainfall layers) and outputs (e.g., inundation maps).
    \item \textbf{The Adapters}: Handle the specific technical details of data conversion (e.g., Rasterio for GeoTIFFs or 
    custom C++ loaders for HIF).
\end{itemize}

This decoupling ensures that the system is highly extensible. For instance, a new "Infiltration Layer" based on a different 
physical formula can be added by simply creating a new adapter, without modifying the core simulation engine.


\subsection{Computational Performance Strategy}

Finally, the strategy addresses the latency problem by prioritizing "Performance by Design". The resolution involves 
implementing the simulation engine in C++, utilizing advanced data structures that map the multilayer cellular automata 
directly onto memory-efficient tensors. By representing the system as a multi-channel grid, the model can leverage 
vectorized operations and parallelism, aiming to achieve simulation speeds that allow for multiple "what-if" scenario 
runs within the decision-making window of a real-world emergency.


\subsection{Summary of the Resolution Approach}

In summary, the proposed methodology does not merely provide a new algorithm but a complete technical framework:

\begin{enumerate}
    \item \textbf{Analytically}, it uses $m:n-CA^k$ to manage environmental complexity. 
    \item \textbf{Architecturally}, it uses Hexagonal Design to ensure flexibility and "Open Model" compliance.
    \item \textbf{Technically}, it utilizes high-performance C++ and the IDRISI format to bridge the gap between GIS 
    and real-time forecasting.
\end{enumerate}

This integrated approach provides the necessary foundation to build a Reactive Digital Twin capable of responding to 
extreme events like the 2024 DANA with the speed and precision required by modern society.








\section{Theoretical Framework: Model Definition $m:n-CA^k$}

The flood simulation model is formalized using an extension of the Cellular Automata (CA) theory, known as $m:n-CA^k$ 
\parencite{Fonseca2011}. This architecture enables the integration of multiple layers of spatial and 
temporal information with different dimensions into a single cellular computational framework.

Furthermore, aligning with recent computational advancements, the model leverages tensorial data structures to provide a 
more accurate, flexible, and computationally efficient framework. This tensor-based approach allows for the 
modeling of continuous spatial domains and continuous cell states within the CA, simplifying the integration of complex 
geographical information \parencite{Fonseca2025}. 

Mathematically, the hydrodynamic cellular automaton is defined as a tuple:

$$A = \langle R, E, V, N, \Psi \rangle$$

Where:
\begin{itemize}
    \item $R$ is the $m$-dimensional cell space.
    \item $E$ is the set of $n$ layers that define the global state and the physical environment of the system.
    \item $V$ represents the vicinity function.
    \item $N$ represents the nucleus function.
    \item $\Psi$ is the combination or evolution function that governs the fluid dynamics.
\end{itemize}


\subsection{Cell Space ($R$) and Layer Definition ($E$)}

The simulation space $R$ is defined on a regular two-dimensional grid $R \subset \mathbb{Z}^2$, where each cell is 
identified by its spatial coordinates $(i, j)$. The physical size of each cell is determined by the spatial resolution 
$\Delta x$.

The system consists of a set $E = \{E_1, E_2, E_3, E_4\}$ of $n=4$ layers. One of the fundamental characteristics of the 
$m:n-CA^k$ model is that the layers may differ in their geometric and temporal dimensions. Through these tensor representations, 
they are defined as follows:

\begin{enumerate}
    \item $E_1$: Topobathymetry ($Z_{topo-bathy}$). A static 2D layer of dimension $X \times Y$. It combines the Digital 
    Elevation Model (DEM) with the bathymetry of the watercourses, providing the absolute elevation of the riverbed: 
    $E_1 \subset \mathbb{R}^{X \times Y}$.
    \item $E_2$: Land Cover (LC). A static 2D layer of size $X \times Y$. It identifies land uses, which are algebraically 
    mapped to Manning roughness coefficients ($\mu_{soil}$) via a state transfer function: 
    $E_2 \subset \mathbb{N}^{X \times Y} \xrightarrow{f} \mathbb{R}^{X \times Y}$.
    \item $E_3$: Water Depth ($WD$). A 2D dynamic layer of dimension $X \times Y$. It represents the main state of the cell 
    at time $t$, storing the water depth in meters: $E_3 \subset \mathbb{R}^{X \times Y}$.
    \item $E_4$: Rainfall ($RF$). A 3D dynamic layer of dimension $X \times Y \times T$, explicitly incorporating the time 
    domain to provide the rainfall rate for each step $\Delta t$: $E_4 \subset \mathbb{R}^{X \times Y \times T}$.
\end{enumerate}


\subsection{Vicinity Function ($V$) and Nucleus Function ($N$)}

The \textbf{Vicinity Function} ($V$) determines the subset of adjacent cells that influence the calculation of the future state of a 
given cell. A Moore neighborhood with radius $r=1$ is adopted, defining a set of 8 adjacent cells connected orthogonally and 
diagonally.

$$V(i,j) = \{ (i+u, j+v) \in R \mid \max(|u|, |v|) \le 1 \}$$

The \textbf{Nucleus Function} ($N$) is projected onto the central cell itself to update layer $E_3$, mapping the spatial 
modifications through local rules: $N(i,j) = \{ (i,j) \}$.


\subsection{Combination / Evolution Function ($\Psi$)}

The evolution function $\Psi: E^t \rightarrow E^{t+1}$ governs the transformation of the system from time $t$ to time $t+1$. 
In this proposed framework, the fluid dynamics computation is divided into two sequential phases: Meteorological Forcing and 
Hydrodynamic Routing.


\subsubsection{Phase 1: Meteorological Forcing}

At the start of iteration $t$, the water volume in the cell is updated by incorporating the contribution from the 3D rainfall 
layer ($E_4$). The intermediate state of the water depth, denoted $WD_{i,j}^*$, is calculated as:

$$WD_{i,j}^* = WD_{i,j}^t + RF(i,j,t)$$


\subsubsection{Phase 2: Kinematics, Viscosity, and Physical Limiters}

Instead of computing momentum conservation via dynamic forces, the evolution core utilizes a multi-physics Kinematic Wave 
approach combined with explicit spatial and temporal volume limiters. This formulation acts as a numerical safeguard against 
instabilities, ensuring accurate solutions when modeling shallow water flows over complex topographies, bottom discontinuities, 
and transient dry-bottom areas \parencite{Aleksyuk2017}.

For each cell in the neighborhood $k \in V(i,j)$, the water surface elevation ($WSE$) is evaluated as $WSE = E_1 + WD^*$. The 
hydraulic gradient $S_k$ towards each valid neighbor ($\Delta WSE_k > 0$) is calculated dividing the elevation drop by the 
directional distance $L_k$ ($\Delta x$ for orthogonal, $\Delta x\sqrt{2}$ for diagonal). The model identifies the maximum local 
slope $S_{max} = \max(S_k)$.

To accommodate fluids with varying rheological properties, the core evaluates the effective fluid velocity ($v_{approx}$) by 
continuously solving two competing flow regimes:

\begin{itemize}
    \item \textbf{Turbulent Regime:} Governs light, low-viscosity fluids (such as pure water). To this end, it is assumed that 
    the flow reaches its terminal velocity almost instantaneously. For the estimation of the kinematic velocity, the model adopts 
    the distributed uniform flow formulation proposed by \parencite{Bates2000} for raster-based flood inundation simulation. 
    Adapting the Manning's equation employed by the authors, and assuming the shallow water approximation where the hydraulic 
    radius is equivalent to the flow depth ($WD^*$), the velocity is defined as a function of the surface roughness coefficient 
    ($\mu_{soil}$) and the local slope gradient ($S_{max}$):
    $$v_{turb} = \frac{1}{\mu_{soil}} \cdot (WD^*)^{2/3} \cdot \sqrt{S_{max}}$$
    
    \item \textbf{Laminar/Viscous Regime:} Dictates the behavior of dense or non-Newtonian flows (such as mud or debris), driven 
    by fluid density ($\rho$) and dynamic viscosity ($\mu_{fluid}$):
    $$v_{lam} = \frac{\rho \cdot g \cdot (WD^*)^2 \cdot S_{max}}{3 \cdot \mu_{fluid}}$$
\end{itemize}

The automaton dictates that the fluid will travel at the lowest of the two velocities, naturally bounding the flow based on its 
dominant physical state:
$$v_{approx} = \min(v_{turb}, v_{lam})$$

Consequently, the theoretical outflow volume ($V_{theor}$) over the time step $\Delta t$ is calculated through the flow 
cross-section, and distributed into directional outflow volumes ($\Delta V_{theor \rightarrow k}$) applying Multiple Flow 
Direction (MFD) weights $W_k$:

$$V_{theor} = v_{approx} \cdot (WD^* \cdot \Delta x) \cdot \Delta t$$
$$\Delta V_{theor \rightarrow k} = V_{theor} \cdot \left( \frac{S_k}{\sum_{n \in V} S_n} \right)$$

To ensure strict mass conservation and prevent the inversion of hydraulic gradients (checkerboard instability), two sequential 
limiters are enforced before state updating:

\begin{enumerate}
    \item \textbf{Spatial Shield (Level Equalization):} The central cell is restricted from dropping its water level beyond half 
    the maximum elevation difference with its lowest neighbor:
    $$\Delta V_{spatial} = \frac{\max(\Delta WSE) \cdot \Delta x^2}{2}$$

    \item \textbf{Temporal Shield (Volumetric CFL Condition):} In traditional Computational Fluid Dynamics (CFD), the 
    Courant-Friedrichs-Lewy (CFL) condition is strictly enforced by constraining the time step ($\Delta t$) or limiting the 
    maximum velocity to ensure that a fluid wave does not traverse a distance greater than the grid resolution:
    $$v \le \frac{\Delta x}{\Delta t}$$
    
    Dynamically altering the global time step in parallel tensor executions is computationally inefficient. Instead, the 
    proposed model enforces the CFL condition volumetrically. By rearranging the classical kinematic condition 
    ($v \cdot \Delta t \le \Delta x$) and multiplying both sides by the flow cross-sectional area ($WD^* \cdot \Delta x$), 
    the inequality is transformed:
    $$(v \cdot \Delta t) \cdot (WD^* \cdot \Delta x) \le \Delta x \cdot (WD^* \cdot \Delta x)$$
    
    This mathematically maps to the model's volumetric routing variables, ensuring that the theoretical intended outflow 
    ($V_{theor}$) does not exceed the absolute total water volume present within the cell ($V_{cell} = WD^* \cdot \Delta x^2$). 
    Consequently, the effectively available volume ($V_{avail}$) is strictly bound by this physical fluid mass and the 
    aforementioned spatial restriction:
    $$V_{avail} = \min(WD^* \cdot \Delta x^2, \Delta V_{theor})$$
\end{enumerate}

A strict scale factor $\alpha$ is derived. If the intended routing volume exceeds the available constraints, all directional flows 
are scaled down proportionally:

$$\alpha = \frac{\min\left(\sum \Delta V_{theor \rightarrow k}, V_{avail}\right)}{\sum \Delta V_{theor \rightarrow k}}$$
$$\Delta V_{out \rightarrow k} = \Delta V_{theor \rightarrow k} \cdot \alpha$$

Finally, the main state variable ($E_3 \equiv WD$) is updated for $t+1$ by translating volumetric transfers ($m^3$) into vertical 
water column variations ($m$). A floating-point precision filter ($\epsilon = 10^{-5}$) is applied to avoid microscopic residual 
accumulations:

$$WD_{i,j}^{t+1} = \max\left(WD_{i,j}^* - \sum_{k \in V} \frac{\Delta V_{out \rightarrow k}}{\Delta x^2} + \sum_{k \in V} \frac{\Delta V_{in \leftarrow k}}{\Delta x^2}, \ 0\right)$$

This integrated tensorial update rigorously preserves volumetric continuity while seamlessly handling the wet-dry interface evolution.











\section{System Macro-Architecture}

The system is designed following a Loose Coupling strategy between two distinct technological domains, effectively 
separating "Data Engineering" (preparation) from "Computational Physics" (simulation). This decoupling is mediated 
by a standardized File-Based Interface, ensuring that the high-level data processing in Python and the high-performance 
execution in C++ remain independent yet perfectly synchronized.

\subsection{The Data-as-Interface Strategy}

Instead of integrating Python scripts directly into the C++ core via interpreters or APIs, the architecture utilizes 
a strictly defined File-Based Interface. The Data Pipeline acts as a "Producer", generating a structured repository 
of geographical and temporal data. The C++ Simulator acts as a "Consumer", ingesting these datasets through a dedicated 
InputAdapter.

The system is a unified environment based on the IDRISI standard. Static layers are stored as IDRISI Raster datasets 
(.doc/.img), while spatiotemporal data is managed through a custom Hierarchical IDRISI Format (HIF).

This strategy provides several key advantages:

\begin{enumerate}
    \item \textbf{Format Homogeneity}: By using a single underlying raster standard (IDRISI), the simulator's input 
    logic is simplified, reducing the overhead of managing multiple third-party libraries like HDF5 in the C++ 
    environment.
    \item \textbf{Technological Specialization}: Python manages complex GIS operations—such as coordinate reference 
    system (CRS) transformations and remote sensing API calls—while the C++ core focuses exclusively on the numerical 
    efficiency of Cellular Automata rules.
    \item \textbf{Immutability and Reproducibility}: Once the pipeline execution is finalized, the input artifacts 
    are fixed. This allows researchers to audit the exact state of the environment before a simulation run, ensuring 
    that results are deterministic and reproducible.
\end{enumerate}

\begin{figure}[htbp]
    \centering
    \includegraphics[width=0.85\textwidth]{chapter_3/system_architecture_overview}
    
    \caption{System architecture overview}
    \label{fig:system_architecture_overview}
\end{figure}








\section{The Hierarchical IDRISI Format (HIF)}

To manage spatiotemporal data (e.g., hourly rainfall or dynamic water levels) without relying on complex binary 
containers like HDF5 or NetCDF, this project introduces the Hierarchical IDRISI Format (HIF). The HIF is a structural 
specification designed to store 3D data cubes using the native 2D IDRISI raster standard as a building block.

\subsection{Conceptual Structure}

The HIF treats time-series data as a discrete sequence of temporal frames. Each frame is stored as an independent 
IDRISI raster pair (.doc for metadata and .img for the data matrix), organized within a hierarchical directory tree. 
This allows for "lazy loading" of simulation steps, where the engine only reads the specific frame required for the 
current time-step $t$.

\subsection{Components of an HIF Dataset}

An HIF dataset consists of three primary elements:

\begin{enumerate}
    \item \textbf{Master Metadata (.hif)}: A global descriptor file that defines the spatial context (extent, resolution, 
    and CRS) shared by all frames in the sequence. It serves as the entry point for the simulator.
    \item \textbf{Temporal Attribute Manifest (attrs.json)}: A specialized metadata file that maps simulation time-steps 
    to specific file paths. It contains the downgrade factors, ISO-8601 timestamps, and the ordered list of frame filenames.
    \item \textbf{Frame Repository}: A flat collection of .img files, each representing a single point in time, without 
    accompanying .doc files.
\end{enumerate}

By adopting this format, the system achieves a balance between the simplicity of ASCII-based/flat-binary rasters and 
the requirement for organized, high-dimensional temporal data.

\subsection{Physical Storage and Directory Topology}

The physical realization of an HIF dataset is organized as a self-contained directory that encapsulates both the global 
spatiotemporal metadata and the individual raster frames. This structure is designed to facilitate modular access: the 
simulation engine can parse the top-level manifest to understand the temporal extent and subsequently perform random 
access to specific time-steps without the need to scan the entire dataset. 

This tree topology illustrates the resulting hierarchy after the Data Pipeline exports a dynamic layer (e.g., rainfall).

\begin{figure}[htbp]
    \centering
    \footnotesize
    \begin{forest}
      for tree={
        folder,   
        font=\ttfamily,
        grow'=0,
      }
    [output/rainfall/
      [\file{rainfall.hif}]
      [\file{attrs.json}]
      [\file{000000.img}]
      [\file{000001.img}]
      [\file{...}]
    ]
    \end{forest}
    \caption{Hierarchical IDRISI Format (HIF) internal directory structure visualized as a tree topology.}
    \label{fig:hif_structure}
\end{figure}











\section{Data Pipeline Design (Pre-processing)}

The core of any flood simulation system lies in its ability to process vast amounts of geospatial 
data accurately and efficiently. The Data Pipeline is designed as a centralized orchestrator 
responsible for the ingestion, transformation, and provision of the physical parameters required 
by the hydraulic model.

The primary objective is to decouple data acquisition from the simulation logic. By doing so, the 
system ensures that changes in data formats or sources do not affect the stability of the 
simulator's core. This framework is built upon three fundamental pillars:

\begin{enumerate}
    \item \textbf{Integrity}: Ensuring that spatial coordinates and resolutions remain consistent 
    across all data layers.
    \item \textbf{Scalability}: Allowing the inclusion of new physical variables (e.g., soil 
    moisture or infiltration) without redesigning the existing architecture.
    \item \textbf{Automation}: Minimizing manual intervention during the preprocessing of raster 
    datasets.
\end{enumerate}

\begin{figure}[htbp]
    \centering
    \includegraphics[width=0.85\textwidth]{chapter_3/high-level_system_context}
    
    \caption{High-Level System Context}
    \label{fig:high-level_system_context}
\end{figure}


\subsection{Data Categorization: Static vs Dynamic Layers}

To optimize processing and memory management, the methodology categorizes geospatial information 
into two distinct types based on their temporal behavior during a simulation:

\begin{itemize}
    \item \textbf{Static Layers}: These represent physical characteristics that remain constant 
    throughout the simulation event. Examples include the Digital Elevation Model (DEM) and Land 
    Use (Roughness). These layers are processed once and serve as the immutable "canvas" of the 
    terrain.
    \item \textbf{Dynamic Layers}: Input data that informs the transition function.These involve 
    variables that change over time, requiring a temporal dimension in their representation. 
    Examples include Rainfall Intensity. The design must account for the sequential reading or 
    generation of these files, ensuring that the simulator receives the correct "snapshot" for 
    each time step.
\end{itemize}

This conceptual separation allows the pipeline to apply different I/O strategies for each type, 
significantly improving the efficiency of the data flow.


\subsection{Architectural Design Patterns}

To achieve a balance between flexibility and rigorous data handling, the system architecture is 
grounded in three fundamental software engineering patterns: Abstraction, Inheritance, and the 
Registry Pattern.


\subsubsection{Abstraction and Encapsulation}

The methodology treats every geospatial data source as a specialized object rather than a simple 
file path or a raw array. By abstracting the complexities of coordinate systems, pixel resolutions, 
and file formats into dedicated classes, the pipeline creates a unified interface. This encapsulation 
ensures that the rest of the simulator does not need to understand the internal structure of the data; 
it only interacts with high-level methods to retrieve physical values.


\subsubsection{Hierarchical Inheritance (Base Classes)}

The design utilizes a hierarchical structure to manage the commonalities and differences between 
data types.

\begin{itemize}
    \item \textbf{Base Components}: A set of abstract base classes defines the "contract" that all 
    specific layers and Input/Output (I/O) handlers must fulfill.
    \item \textbf{Specialization}: This allows the system to distinguish between the static nature 
    of a Digital Elevation Model (DEM) and the time-varying nature of rainfall intensity, while 
    still sharing the same underlying logic for spatial referencing and memory management.
\end{itemize}

\begin{figure}[htbp]
    \centering
    \includegraphics[width=0.85\textwidth]{chapter_3/data_generator_hierarchy}
    
    \caption{Data Generator Hierarchy}
    \label{fig:data_generator_hierarchy}
\end{figure}


\subsubsection{The Registry Pattern for Extensibility}

One of the most critical requirements for the TFM is the ability to integrate new physical variables 
(layers) without modifying the core orchestrator. To solve this, the methodology adopts a 
\textbf{Registry Pattern}.

In this approach, the Data Pipeline acts as a central hub that does not "know" about specific layers 
(like Rainfall or Roughness) at compile-time. Instead, each new layer "registers" itself into a 
dictionary-like structure. When the simulation starts, the pipeline simply iterates through this 
registry to initialize the required components. This design achieves Low Coupling, as adding a fifth 
or sixth layer only requires creating a new class that follows the established base contract, leaving 
the main execution logic untouched.


\subsubsection{Decoupling I/O from Logic}

Finally, the architecture enforces a strict separation between Data Access (reading/writing files) and 
Data Logic (processing the raster values). By isolating I/O operations into dedicated I/O Handlers, 
the system becomes agnostic to the storage medium. If the project were to move from local IDRISI files 
to a cloud-based database or a different file format (like NetCDF), only the I/O handler would need to 
be updated, while the layer generation logic remains intact.


\subsection{Data Flow Logic}

The movement of data through the system is governed by a deterministic lifecycle, ensuring that the 
simulation engine receives a spatially and temporally synchronized "snapshot" of the environment. This 
flow is divided into four primary stages:


\subsubsection{Configuration-Driven Discovery}

The process begins with the ingestion of a centralized configuration manifest (typically in YAML format), 
which serves as the "Control Plane" for the entire data pipeline. Unlike a discovery-based system, the pipeline 
maintains a pre-defined registry of core physical layers (e.g., Elevation, Rainfall, Roughness) essential for 
the simulation. The configuration file is utilized to parameterize these layers and define the global execution 
context.

This stage manages three critical data dimensions:

\begin{enumerate}
    \item \textbf{Global Simulation Metadata}: It establishes the temporal boundaries of the scenario (start and 
    end timestamps) and the workspace topology (output directory and naming conventions). This ensures that all 
    generated layers are synchronized within the same spatiotemporal window.
    \item \textbf{Layer Parameterization}: For each architectural layer, the configuration provides specific 
    execution parameters. This includes spatial resolutions, interpolation models (e.g., Variogram types for 
    rainfall), and geographical constraints (bounding boxes). This allows the same pipeline logic to be reused 
    across different study areas by simply adjusting the numerical constants in the config file.
    \item \textbf{Visualization Manifest}: The pipeline identifies which layers require graphical rendering for 
    research purposes. By defining a "visualize" list in the configuration, the system selectively triggers the 
    visualization engine post-generation, optimizing execution time by avoiding unnecessary plots for background 
    layers.
\end{enumerate}

By separating the layer definitions (hard-coded in the pipeline's logic for structural integrity) from their 
operational parameters (externalized in the configuration), the system achieves a high degree of reproducibility 
and scientific flexibility.

\subsubsection{Spatial Synchronization and Metadata Validation}

Before any numerical processing occurs, the pipeline performs a Spatial Integrity Check. Flood simulations 
are highly sensitive to spatial misalignment. The flow logic enforces a "Master Grid" approach, where 
secondary layers (such as rainfall or roughness) must be validated against the primary Elevation Layer (DEM).

\begin{itemize}
    \item \textbf{Coordinate Reference System (CRS) Alignment}: Ensuring all data layers use the same 
    projection.
    \item \textbf{Resolution Matching}: Verifying that pixel sizes are consistent across all rasters to avoid 
    interpolation errors during the simulation.
    \item \textbf{Extent Clipping}: Ensuring that all dynamic inputs cover the exact bounding box of the 
    static terrain.
\end{itemize}


\subsubsection{Sequential Execution: The Static-Dynamic Split}

The data flow logic branches based on the temporal nature of the data:

\begin{enumerate}
    \item \textbf{The Static Phase (Setup)}: The pipeline triggers the generation and storage of immutable 
    layers (Elevation and Roughness). These are processed once and held in a "Ready-to-Simulate" state.
    \item \textbf{The Dynamic Loop (Runtime)}: For variables like Rainfall, the pipeline enters a cyclic flow. 
    For every simulation time step, the logic dictates a "Fetch-Transform-Provision" cycle, ensuring that the 
    simulator always has access to the most recent temporal state without overloading the system's memory.
\end{enumerate}


\subsection{Validation and Logging Strategy}

The robustness of the data pipeline is sustained by a dual-layer strategy: rigorous pre-execution validation 
to ensure structural and spatial coherence, and comprehensive logging to provide full transparency during the 
transformation processes.


\subsubsection{Structural and Dependency Validation}

Before processing any raster data, the system must validate the logical structure of the simulation request. 
This is approached using Graph Theory principles.

\begin{itemize}
    \item \textbf{Directed Acyclic Graph (DAG)}: The dependencies between layers (e.g., Rainfall requiring an 
    Elevation model for clipping) are treated as nodes and edges in a graph. The methodology requires a 
    topological sort to ensure that no circular dependencies exist and that every component is initialized 
    only when its requirements are met.
    \item \textbf{Schema Enforcement}: The configuration files are validated against a predefined schema to 
    ensure that essential parameters, such as coordinate reference systems or simulation timeframes, are 
    present and logically consistent (e.g., start\_date must precede end\_date).
\end{itemize}

\begin{figure}[htbp]
    \centering
    \includegraphics[width=0.7\textwidth]{chapter_3/layer_dependencies}
    
    \caption{Layer Dependencies}
    \label{fig:layer_dependencies}
\end{figure}


\subsubsection{Spatial Integrity and Physical Constraints}

The core of the validation strategy lies in ensuring "Spatial Coherence". Since the hydraulic model operates 
on a numerical grid, the pipeline enforces strict spatial constraints:

\begin{itemize}
    \item \textbf{Domain Consistency}: Validating that all input datasets overlap with the requested study area.
    \item \textbf{Resolution Uniformity}: Detecting and flagging discrepancies in pixel size that could lead 
    to significant calculation errors in the fluid dynamics solver.
    \item \textbf{No-Data Handling}: A theoretical framework for "Nodata" values is established, ensuring that 
    missing geographic information does not propagate as zero-values, which would lead to physical inaccuracies 
    in the flood extent.
\end{itemize}


\subsubsection{Observability through Multi-Sink Logging}

In long-running environmental simulations, observability is critical for both real-time monitoring and post-mortem 
analysis. The methodology adopts a Multi-Sink Logging Pattern:

\begin{enumerate}
    \item \textbf{Transient Output (Standard Stream)}: Focused on high-level operational status (Success, Info, Error). 
    It provides the user with immediate feedback on the pipeline's progress without cluttering the interface 
    with technical debt.
    \item \textbf{Persistent Output (File-based Log)}: A detailed, low-level record of every transformation, I/O 
    operation, and minor warning. This serves as a "black box" for researchers to audit the data generation process.
    \item \textbf{Severity Layering}: Messages are categorized by intensity (Trace, Debug, Info, Success, Warning, 
    Error). This allows for granular filtering: while a researcher might only want to see successful milestones, 
    a developer can access the "Trace" level to inspect the exact metadata of a IDRISI being processed.
\end{enumerate}

\begin{figure}[htbp]
    \centering
    \includegraphics[width=0.7\textwidth]{chapter_3/logging_architecture}
    
    \caption{Logging Architecture}
    \label{fig:logging_architecture}
\end{figure}








\section{Simulation Engine Design}

\subsection{The Simulator Core: State inference logic}

\subsection{Hexagonal Architecture: Decoupling Using Ports and Adapters}

\subsection{Advantages of the design for physical layer extensibility}

\subsection{Data Structure and Tensor Model}

\subsection{Representation of the system as a multi-channel tensor}

\subsection{Layer Metadata Specification}
